\documentclass[a4paper,12pt]{article}
\usepackage[left=2.5cm,top=3cm,right=2.5cm,bottom=3cm]{geometry}
\usepackage[english]{babel}
\usepackage[utf8]{inputenc}
\usepackage{amssymb}
\usepackage{amsmath}
% ========= Fix for long URLs breaking out of page =========
\PassOptionsToPackage{hyphens}{url}
\usepackage{url}
\usepackage[hyphens]{hyperref}
\hypersetup{breaklinks=true}
%Figures
\usepackage{graphicx}
\usepackage{caption}
\usepackage{subcaption}
\usepackage[table]{xcolor}
\usepackage{longtable}
%Tables
\usepackage{booktabs}
%Bold in tables
\newcommand{\bftab}{\fontseries{b}\selectfont}
\newsavebox\CBox
\def\textBF#1{\sbox\CBox{#1}\resizebox{\wd\CBox}{\ht\CBox}{\textbf{#1}}}
%dashed line
\usepackage{array}
\usepackage{arydshln}
\setlength\dashlinedash{0.2pt}
\setlength\dashlinegap{1.5pt}
\setlength\arrayrulewidth{0.3pt}
% symbols
\newcommand{\Lagr}{\mathcal{L}}
\usepackage{gensymb} % degree
%Includes "References" in the table of contents
\usepackage[nottoc]{tocbibind}
% Harvard citation style
\usepackage{natbib}
% KM comment macro
\newcommand{\km}[1]{
{\color{green}{KM: \it{[#1]}}}
}
\usepackage{multicol}
\begin{document}
\begin{titlepage}
    \begin{center}
        % Imperial logo (put the correct file name here)
        \includegraphics[width=0.45\textwidth]{imperial_logo.png}\\[2cm]
        {\Large \textsc{Laboratory in Applied Machine Learning}}\\[0.6cm]
        {\large \textsc{Imperial College London}}\\[0.3cm]
        {\large \textsc{Department of Electrical and Electronic Engineering}}\\[1.2cm]
        \rule{\textwidth}{0.4pt}\\[0.5cm]
        {\LARGE \textbf{sEMG-Based Smart Silent Speech Reader (SSSR)}}\\[0.3cm]
        {\large Final Report}\\[0.5cm]
        \rule{\textwidth}{0.4pt}\\[2cm]
    \end{center}
    \vspace{0.5cm}
    \begin{minipage}{0.45\textwidth}
        \begin{flushleft} \large
            \textit{Authors:}\\
            Yash M Agarwal\\
            Ivan Plakhotnyk\\
            Dorijan Donaj Maga\v{s}i\'{c}
        \end{flushleft}
    \end{minipage}
    \hfill
    \begin{minipage}{0.45\textwidth}
        \begin{flushright} \large
            \textit{Supervisors:}\\
            Prof.\ Krystian Mikolajczyk\\
            Dr.\ Yassir Jedra\\
            Dr.\ Abd Al Rahman Abu Ebayyeh\\
        \end{flushright}
    \end{minipage}
    \vfill
    \begin{center}
        {\large March 2026}
    \end{center}
\end{titlepage}
\begin{abstract}
Silent speech interfaces (SSIs) enable communication by decoding facial muscle activity during unvoiced articulation, offering a lifeline for individuals unable to vocalise. However, existing systems typically rely on research-grade electrode arrays that are prohibitively expensive for practical deployment. We present the Smart Silent Speech Reader (SSSR), a consumer-grade sEMG-to-text system built around four MyoWare~2.0 sensors and an Arduino Uno~R4 at approximately one-fifth the cost of the Gaddy \& Klein benchmark setup. Our pipeline pairs a lightweight convolutional--recurrent adaptor (1.8--10.1\,M parameters) with a frozen large language model (LLM) to decode unvoiced EMG into text. A systematic comparison of eight frozen LLMs reveals that Gemma-2-2b-it achieves 0.3285 WER on the Gaddy closed-vocabulary benchmark; a 33\% relative improvement over the prior state of the art (0.49 WER with LLaMA-3.2-3B). A nine-stage ablation study across four LLMs identifies prompt-consistent inference, output layer normalisation, and dropout as universally beneficial, while showing that the optimal training recipe is LLM-specific. We further demonstrate that performance degrades approximately linearly from eight to three EMG channels, with even our three-channel configuration outperforming the prior eight-channel baseline. Finally, on a custom 10-word phonetically distinct vocabulary, three channels and 250 training samples suffice to reach 92\% word accuracy. These results show that affordable, low-channel-count hardware combined with an appropriate frozen LLM can deliver competitive silent speech recognition without research-grade equipment.
\end{abstract}

\pagebreak
\tableofcontents
\pagebreak
\section{Introduction}

Silent Speech Interfaces are designed to recognise speech without sound, thereby allowing communication in scenarios where vocalisation is impossible or undesirable. sEMG is appropriate for this task, since it captures muscle activity from facial articulators even during silent mouthing.

The Smart Silent Speech Reader (SSSR) project is focused on the use of machine learning to assign meaning to unvoiced sEMG signals in terms of word representations. The previous Imperial project introduced a small-vocabulary Arduino-based classifier \citep{sun2022}, but performance there was largely limited by noise, user dependence, and constrained expressiveness.

Our project develops this capability further by investigating better sensing and modelling approaches that could support a larger vocabulary. This aligns with recent work showing that LLM-based EMG-to-text pipelines can interpret unvoiced EMG with strong data efficiency \citep{mohapatra2025} and that multimodal LLM-enhanced models can significantly improve silent speech recognition \citep{benster2024}. This involves overcoming issues such as weak and variable EMG signals, missing voiced data, and microcontroller hardware constraints. The overall result is a practical, demonstrable system for silent communication based on unvoiced EMG. This system ultimately aims to assist individuals who cannot vocalise speech (e.g., due to neurological or muscular impairment) and to enable silent or private communication in everyday environments.

\subsection{Aims \& Objectives}

We aim to build a system that can convert EMG signals captured from targeted facial muscles into meaningful word representations. To achieve this, we propose to train a deep learning-based pipeline on unvoiced EMG signals, using textual transcriptions as ground-truth labels during training. At inference time, the model should operate solely on unvoiced EMG, enabling communication for scenarios where the user cannot produce sound but needs to convey information.

The objectives of this work include:
\begin{itemize}
    \item To design a sensing setup targeted at facial muscles involved in articulation.
    \item To collect a synchronised dataset of unvoiced EMG signals aligned with textual labels.
    \item To develop signal processing and feature extraction methods suitable for EMG and auxiliary sensors.
    \item To train and evaluate deep learning models capable of mapping EMG sequences to words or sentences.
    \item To explore extended vocabulary recognition, beyond the small sets (classification task) used in previous implementations.
    \item To deploy a simplified version suitable for real-time demonstration using Arduino and a computer (for the deep learing models).
\end{itemize}

\section{Literature Review}

\subsection{2022 AML Closed-Vocabulary SSR}

The 2022 Imperial College AML Lab project demonstrated that a functional silent speech interface can be built using inexpensive surface EMG sensors and an Arduino microcontroller \citep{sun2022}. Their system used three MyoWare sensors placed on key facial muscles and applied classical preprocessing such as wavelet denoising, transient-state removal, and sliding-window feature extraction. A range of machine learning models were tested, from Random Forests, SVMs, and KNN to CNN architectures, with accuracies varying from 76 to 86 percent depending on the ML model. Upon testing new subjects, accuracy sharply dropped to 31.55\%.

\subsection{Foundational EMG-to-Speech Work}

The foundational benchmark for modern silent speech recognition was established by \citet{gaddy2021}, who developed an end-to-end system for converting facial EMG signals from eight electrodes into speech audio. Their approach predicts mel-frequency cepstral coefficients (MFCCs) from raw EMG and synthesises audio using a WaveNet vocoder. Two key architectural contributions drove their improvements: replacing hand-designed EMG features with learned convolutional features extracted directly from raw signals, and substituting LSTMs with a Transformer encoder using relative position embeddings. An auxiliary phoneme prediction loss provided additional supervision during training. To handle the temporal misalignment between silent and vocalised speech, the model uses dynamic time warping (DTW) to align silent EMG recordings against their vocalised counterparts during training. With approximately 40 million parameters trained on 19 hours of single-speaker data, the system achieved a word error rate (WER) of 42.2\% on open-vocabulary transcription, a substantial reduction from the previous 68.0\%. This work defined the dataset, electrode placement protocol, and evaluation methodology that most subsequent studies build upon.

\subsection{LLM-Augmented Silent Speech Recognition}

The current state of the art in unvoiced silent speech recognition on the Gaddy and Klein benchmark is the 2025 study \textit{``Can LLMs Understand Unvoiced Speech?''} \citep{mohapatra2025}. This work proposes a two-stage EMG-to-LLM pipeline in which an adaptor network maps high-frequency EMG into the input embedding space of a frozen large language model (LLaMA-2-7B and LLaMA-3.2-3B). The adaptor is composed of 1D convolutional layers for downsampling the high-rate EMG signal ($\leq$800 Hz), residual blocks for temporal feature extraction, BiLSTM layers for modelling sequential dynamics, and fully connected layers that project the extracted features into the embedding dimension of the target LLM. The resulting sequence of EMG embeddings is concatenated with text prompts and fed into the frozen LLM for autoregressive decoding.

Evaluated on the closed-vocabulary subset of the Gaddy and Klein dataset (67 words, approximately 26 minutes of data), the architecture achieves a WER of 0.49 with only 6 million trainable parameters, significantly outperforming application-specific EMG models with 54 million parameters that reach only 0.75 WER. The study further demonstrates that handcrafted EMG features improve LLM performance over raw signals, and that LLMs learn more effectively from audio than from EMG, highlighting the inherent difficulty of the unvoiced modality. An additional finding is that person identification from unvoiced EMG achieves 96\% accuracy, which motivates the development of personalised rather than speaker-independent models.

\subsection{Further Developments in Silent Speech Research}

Beyond the core works discussed above, several parallel developments have shaped the current landscape of silent speech research. \citet{benster2024} introduced MONA LISA, a cross-modal framework that aligns EMG and audio representations through contrastive learning objectives (crossCon and supTcon), performing temporal alignment in learned latent space rather than on raw signals. Combined with LLM-based rescoring of beam search outputs at inference time, the system achieves 12.2\% WER on silent EMG---the first to break the 15\% threshold for non-invasive silent speech---and 3.7\% on vocal EMG. In a different linguistic context, \citet{song2023} demonstrated that Transformer decoders substantially outperform LSTM decoders for syllable-level sequential decoding of Mandarin Chinese from 64-channel high-density sEMG, achieving 5.14\% CER and 96.37\% phrase accuracy. \citet{xie2025} further extended the field to Chinese with the first end-to-end silent EMG-to-text system, using a Transformer encoder with CTC decoding, auxiliary pinyin generation, and adversarial session classification to achieve 38.0\% CER across five recording sessions. Their ablation study confirms that facial electrode channels contribute significantly more than laryngeal channels in silent mode, consistent with the intuition that silent speech involves mouthing without vocal fold vibration.

On the practical side, \citet{inoue2025} address the problem of heterogeneous electrode configurations by proposing tokeniser strategies that allow a single HTNet+Conformer model to handle varying sensor setups, achieving 95.3\% word classification accuracy on a 64-word vocabulary and demonstrating cross-language transfer from Japanese to English. \citet{hwang2026} complement this with a systematic analysis of channel efficiency on the Gaddy dataset, showing that optimal reduced-channel subsets exploit complementary spatial coverage rather than simply retaining the individually strongest channels. Their channel dropout pre-training strategy, in which channels are randomly masked during full eight-channel training before fine-tuning on a target subset, consistently outperforms training from scratch in four-to-six channel configurations---a finding directly relevant to our project, which operates with fewer electrodes than the Gaddy benchmark. The emergence of self-supervised foundation models further advances the field: \citet{fasulo2026} introduce TinyMyo, a 3.6M-parameter Transformer encoder pre-trained via masked reconstruction that achieves competitive performance across gesture classification, kinematic regression, and silent speech tasks (33.54\% WER on speech synthesis) while being the first EMG foundation model deployed on an ultra-low-power microcontroller. Finally, the paradigm of projecting biosignal embeddings into frozen LLM token spaces extends beyond EMG: \citet{mishra2025} demonstrate with Thought2Text that EEG signals evoked by visual stimuli can be decoded into text using instruction-tuned LLMs fine-tuned through a CLIP-aligned multi-stage pipeline, validating the generality of the adaptor-based approach used by \citet{mohapatra2025} across different biosignal modalities.

\section{Planning \& Collaboration}

The development of the Smart Silent Speech Reader (SSSR) system was divided into two parallel workstreams---Hardware Setup and Software Setup---whose timeline is shown in Figure~\ref{fig:gantt}. The project ran throughout the Spring term (January--March 2026), with an additional two-week buffer before the demonstration to repeat experiments, retrain models, or resolve hardware issues.

\begin{figure}[ht]
    \centering
    \includegraphics[width=\textwidth]{figures/image1.png}
    \caption{Gantt chart outlining the planned development schedule for the SSSR system.}
    \label{fig:gantt}
\end{figure}

\subsection{Hardware Setup (Weeks 1--5)}

This workstream established a reliable multi-channel EMG acquisition system using the Arduino Uno R4 and MyoWare 2.0 sensors.

\begin{itemize}
    \item \textbf{Component preparation (Week 1).} Hardware was purchased and the UNO R4 platform was explored to understand its ADC and DMA capabilities.
    \item \textbf{Sampling and transfer tests (Weeks 2--3).} We verify stable sampling and USB data transfer.
    \item \textbf{System assembly (Week 4).} Sensors were mounted and tested.
    \item \textbf{Experiment protocol (Weeks 5).} A standardised procedure for electrode placement, repetitions, session structure were finalised.
    \item \textbf{First data collection (Weeks 5).} Initial EMG recordings of isolated words and phrases were gathered for quality assessment and model development.
\end{itemize}

\subsection{Software Setup (Weeks 3--10)}

Software development proceeded in parallel with hardware work.

\begin{itemize}
    \item \textbf{LLM familiarisation (Weeks 3-4).} We studied the use of input embeddings, tokenisers, and LLaMA frozen-model.
    \item \textbf{Data quality assessment (Weeks 5--6).} Recorded EMG was analysed and we decided if we need to repeat the experiment.
    \item \textbf{Adaptor implementation (Weeks 6).} Baseline adaptors were built to project EMG into the LLM embedding space.
    \item \textbf{Adaptor--LLM integration (Weeks 7).} EMG embeddings were combined with prompt embeddings and validated.
    \item \textbf{Training and evaluation (Weeks 8).} The adaptor was trained using cross-entropy loss and was evaluated.
    \item \textbf{Revision and adjustment (Weeks 9).} Depending on results, we would have recollected the data, or refined preprocessing.
\end{itemize}

\subsection{Final Integration and Demonstration (Weeks 9--10)}

The last phase focused on real-time pipeline implementation and preparing for the demo. We were going to try to make systems that are scalable and easy to interface with. With this design, if the results are unsatisfactory, we would have been able to add more sensors and collect more data without significant changes to the setup. We also explored other methods that we might discover in the future. We have also completed the final project report during this period, writing it in parallel with the system development to ensure that documentation evolves alongside our technical progress.

\section{Hardware}
\label{sec:hardware}

\subsection{sEMG Signal Characteristics}

sEMG measures the voltage between two electrodes placed on the skin above a muscle during contraction. The recorded signal does not arise from a single source but from the superposition of many individual motor unit action potentials within the muscle. This is a particular issue for facial muscles because of their overlapping structure, which results in crosstalk between adjacent muscles and can degrade data quality. To mitigate this, the sensor must be placed close to the middle of the muscle of interest (the muscle belly) and kept at a consistent location between experiments.

sEMG is a weak signal---up to 1\,mV peak-to-peak for facial muscles, compared with up to 10\,mV for large limb muscles such as the biceps. Bringing the signal into the Arduino sensing range (0--5\,V) requires significant amplification; however, noise is amplified alongside the signal. Bandpass filtering (20--500\,Hz, where most of the EMG energy lies) is the most common pre-processing technique to suppress this intrinsic noise. We set the sampling frequency to 1,000\,Hz, satisfying the Nyquist criterion for the 500\,Hz upper band edge. Additionally, AC mains interference at 50\,Hz falls within the signal band and can be substantial; a notch filter at 50\,Hz is therefore applied as a further processing step.

\subsection{Hardware Selection}

We chose the Arduino Uno R4, a significant upgrade over the R3. It offers a 16$\times$ more precise ADC (14-bit vs.\ 10-bit), faster conversions, and configurable DMA for stable data transfer. Although the R4 has a smaller community than the R3, the performance advantages greatly outweigh this drawback.

We selected four MyoWare 2.0 sensors for EMG acquisition. The first-generation MyoWare sensors were successfully used by previous students on the same lab project \citep{sun2022}, establishing confidence in the platform. MyoWare also offers an Arduino-compatible extension board (\pounds9) supporting up to six sensors, which simplifies the setup and reduces the risks associated with custom hardware. The sensors are specifically designed for Arduino, with a fixed 200$\times$ amplification gain tuned for the application. We used four ArduLink sensor extension boards and four short 3.5\,mm jack-to-jack cables to connect the sensors to the shield.

\subsection{Data Acquisition and Firmware}

The system's firmware manages the real-time pipeline between the muscle sensors and the host computer. As illustrated in Figure~\ref{fig:data_pipeline}, the process is divided into two primary stages: acquisition and transfer.

To maintain a stable 1 kHz sampling rate across four channels, we utilize a hardware timer to trigger an interrupt-driven ADC Scan Mode. This ensures ``simultaneous'' sampling of all facial muscles. The resulting data is stored in a ring buffer on the Arduino Uno R4. To optimize throughput, data is only transmitted once the accumulated samples match the size of the UART TX buffer. This non-blocking serial transfer prevents data loss and minimizes CPU overhead.

Finally, the host-side software receives the raw binary data, parses it into arrays, and saves the recording. A simple character-based protocol (sending `S' for Start or `E' for End) allows the host to control the Arduino's recording state remotely.

\begin{figure}[ht]
    \centering
    \includegraphics[width=\textwidth]{figures/image2.png}
    \caption{Data pipeline architecture. Acquisition (bottom to center) captures the filtered 0--5V analog muscle signals into a ring buffer via a 1 kHz timer interrupt; Transfer (center to top) manages the non-blocking serial relay.}
    \label{fig:data_pipeline}
\end{figure}

Table~\ref{tab:hardware_comparison} compares our setup with the research-grade system used by \citet{gaddy2021}, whose dataset serves as the primary benchmark for our ML experiments. With a strict project budget of \pounds300, we achieved a 5$\times$ cost reduction compared to the Gaddy benchmark. Gaddy's monopolar electrode configuration picks up significant common-mode noise across the entire head, whereas our bipolar (differential) configuration rejects this noise at the muscle site before it reaches the Arduino. This allows us to achieve a high signal-to-noise ratio with only a 14-bit ADC and 4 channels, in place of a 24-bit ADC and 8 channels. Our setup produces a scientifically valid dataset at approximately 20\% of the cost.

\begin{table}[ht]
\centering
\caption{Comparison of our hardware setup with the Gaddy \& Klein (2020) benchmark.}
\label{tab:hardware_comparison}
\begin{tabular}{p{3.5cm}p{5.5cm}p{5.5cm}}
\toprule
\textbf{Parameter} & \textbf{Gaddy \& Klein (2020)} & \textbf{Our Setup} \\
\midrule
Acquisition Board & OpenBCI Cyton Biosensing Board (research-grade) & Arduino Uno R4 + MyoWare Arduino Shield \\
ADC Resolution & \textbf{24-bit} & \textbf{14-bit} \\
Channels & 8 EMG channels & 4 EMG channels \\
Electrode Config. & \textbf{Monopolar} (shared reference behind ear) & \textbf{Bipolar} / differential \\
Sample Rate & 1,000 Hz & 1,000 Hz \\
Bandpass Filtering & \textbf{No hardware filtering} & 20--500 Hz first-order bandpass \\
Electrodes & \textbf{Gold-plated} reusable electrodes with Ten20 conductive paste & Disposable \textbf{Ag/AgCl (silver)} snap electrodes \\
Cost & $\sim$\pounds1,000+ & $\sim$\pounds200 \\
\bottomrule
\end{tabular}
\end{table}

To align with the high-variance features identified by \citet{gaddy2020}, we targeted four muscle groups critical for articulatory movement: the Orbicularis oris (lips), Masseter (jaw), Mentalis (chin), and Zygomaticus major (upper mouth corner). Focusing on the perioral and chin regions concentrated our sensors on the sites of maximum displacement during speech.

\begin{figure}[ht]
    \centering
    \begin{subfigure}[b]{0.4\textwidth}
        \centering
        \includegraphics[width=\textwidth]{figures/image3.jpeg}
    \end{subfigure}
    \hfill
    \begin{subfigure}[b]{0.45\textwidth}
        \centering
        \includegraphics[width=\textwidth]{figures/image4.jpeg}
    \end{subfigure}
    \caption{Target muscles and sensor placement.}
    \label{fig:target_muscles}
\end{figure}

\section{Data Collection Interface}

To streamline the recording process and ensure consistency across sessions, we developed a portable HTML-based interface for phrase generation and presentation. The core workflow is simple: the participant places the sEMG sensors on their face and then silently mouths the words or phrases displayed on screen, without vocalising. The interface generates script-driven random prompts from a closed vocabulary and presents them one at a time, advancing either automatically on a timer or manually at the participant's pace.

The session configuration panel exposes several adjustable parameters, including the total number of utterances, a countdown delay before each prompt, the mouthing duration (how long the phrase stays on screen), and the rest interval between prompts. A start index field allows a session to be resumed partway through if recording is interrupted. The interface also supports keyboard shortcuts for pausing, repeating, skipping, and ending a session, which proved useful for maintaining a natural recording flow without needing to reach for on-screen controls.

\begin{figure}[ht]
    \centering
    \includegraphics[width=0.85\textwidth]{figures/image5.png}
    \caption{Data Collection Interface, showing several adjustable parameters and phrase generator screen.}
    \label{fig:data_collection_interface}
\end{figure}

We conducted two separate data collection trials, each targeting a different vocabulary design. The first trial followed the closed-vocabulary date-and-time paradigm used by \citet{gaddy2021}, drawing from 67 unique words spanning times (e.g.\ ``08:33 PM''), weekday-dates (e.g.\ ``Friday June 11''), and date-years (e.g.\ ``September 12 2013''). This produced four phrase types --- times with and without AM/PM markers, weekday-month-date combinations, and month-date-year combinations --- recorded across two sessions totalling 438 utterances, with a roughly balanced split of 236 and 202 utterances per session. The phrase distribution was weighted toward weekday-month-date combinations (227 of 438), reflecting the larger number of possible permutations in that category.

The second trial shifted to a single-word vocabulary of 10 phonetically diverse fruit and vegetable names --- pomegranate, dragonfruit, mangosteen, guava, cranberry, sprouts, artichoke, cauliflower, asparagus, and zucchini --- chosen to maximise articulatory difference across the sEMG channels. This trial ran for three sessions producing 757 utterances in total (251, 250, and 256 per session), with all 10 words uttered approximately equally at around 76 repetitions each. The balanced distribution was enforced by the script's randomisation logic to ensure each word class received sufficient representation for downstream classification.

In both trials, each per-phrase EMG recording was saved as a separate CSV file, preserving the alignment between the displayed prompt and the corresponding muscle activity signal. This structure simplifies later preprocessing, as each file maps directly to a single labelled utterance.

\section{Software Pipeline}
\label{sec:software_pipeline}

The software pipeline consists of the following three stages:

\begin{enumerate}
    \item Stage 1: Data processing \& feature extraction from raw EMG signals
    \item Stage 2: Training an adaptor - a specialized neural network to learn mapping between EMG signals and LLM input embedding space
    \item Stage 3: Using a frozen LLM to predict words
\end{enumerate}

This pipeline is also depicted in Figure~\ref{fig:pipeline_overview}.

\begin{figure}[ht]
    \centering
    \includegraphics[width=\textwidth]{figures/image6.png}
    \caption{Pipeline overview.}
    \label{fig:pipeline_overview}
\end{figure}

\subsection{Stage 1: Data Processing \& Feature Extraction}

The collected EMG signals through the Arduino and MyoWare 2.0 sensors (4 channels) are recorded as individual CSV files, one per utterance. Each file contains raw ADC readings at 1\,kHz across all channels. The signal conditioning steps described in Section~\ref{sec:hardware}---bandpass filtering (20--500\,Hz) and 50\,Hz notch filtering---are applied in hardware by the MyoWare sensors prior to digitisation.

From the filtered signal, we extract 14 handcrafted (HC) features per channel per time step: 5 time-domain features and 9 frequency-domain features derived from short-time Fourier transform (STFT) magnitude bins. This feature set follows the recommendations of \citet{gaddy2020} and \citet{mohapatra2025}, who demonstrate that handcrafted features consistently outperform raw EMG as input to LLM-based adaptors. Features are computed over short overlapping windows and concatenated across all channels, producing a feature vector of dimension $C \times 14$ at each time step (e.g., 112 for 8 channels, 56 for 4 channels).

For the Gaddy dataset, we use pre-extracted features provided by \citet{mohapatra2025}. For our personal dataset, we apply the same feature extraction pipeline to our raw recordings. The extracted features are then normalised to bring their distribution in line with the Gaddy feature range. Finally, the normalised feature sequences are saved in the format expected by the adaptor's input layer (Section~\ref{sec:adaptor}).

\subsection{Stage 2: EMG Adaptor}
\label{sec:adaptor}

This module consists of a convolutional layer, three ResBlocks, a convolutional layer, a BiLSTM layer, a convolutional layer and finally a linear layer. We define a parameter model\_size which controls the adaptor's capacity, or its width. Our setup consists of 4 channels, but as an example, for the Gaddy dataset, if we consider 8 channels, and 14 handcrafted features, then the input dimension to the adaptor is a multiplication of the channels and features, 112 here. An example flow of model\_size as 768 and 384 is shown in Table~\ref{tab:adaptor_flow}. It is also worth noting that the data is downsampled by a factor of 48 from $T$ to $T/48$ in the last layer of the adaptor, which is then linearly mapped to the input embedding space dimension of the chosen LLM (depicted as llm\_hidden). Every LLM has a different embedding space, and we depict a few examples in Table~\ref{tab:llm_hidden} that illustrate the interplay between model\_size and llm\_hidden. The model\_size parameter is crucial to control overfitting or underfitting as we experiment with different values. Notably, our adaptor architecture is inspired by \citet{mohapatra2025}.

\begin{table}[ht]
\centering
\caption{Example data flow through the adaptor for model\_size 768 and 384.}
\label{tab:adaptor_flow}
\label{tab:adaptor_flow_384}
\begin{tabular}{lcc}
\toprule
\textbf{Layer} & \textbf{MS = 768} & \textbf{MS = 384} \\
\midrule
ResBlock 1 & $112 \to 768$ & $112 \to 384$ \\
ResBlock 2 & $768 \to 384$ & $384 \to 192$ \\
ResBlock 3 & $384 \to 192$ & $192 \to 96$ \\
BiLSTM     & $192 \to 384$ & $96 \to 192$ \\
Linear     & $384 \to \text{llm\_hidden}$ & $192 \to \text{llm\_hidden}$ \\
\bottomrule
\end{tabular}
\end{table}

\begin{table}[ht]
\centering
\caption{Interplay between model\_size and llm\_hidden for different LLMs.}
\label{tab:llm_hidden}
\begin{tabular}{clll}
\toprule
\textbf{model\_size} & \textbf{LLM} & \textbf{Linear layer} & \textbf{Total params} \\
\midrule
768  & Gemma (hidden=2304)   & $384 \to 2304$ & $\sim$6.0M \\
768  & Qwen (hidden=2048)    & $384 \to 2048$ & $\sim$5.9M \\
768  & LLaMA-3.2-3B (hidden=3072)& $384 \to 3072$ & $\sim$6.3M \\
384  & Gemma (hidden=2304)   & $192 \to 2304$ & $\sim$1.9M \\
1024 & Gemma (hidden=2304)   & $512 \to 2304$ & $\sim$10.1M \\
\bottomrule
\end{tabular}
\end{table}

We broadly perform three experiments on the adaptor stage:

\begin{enumerate}
    \item Grid search: Varying model\_size across 2 LLMs (and a comparison between 8 LLMs for fixed model\_size)
    \item Ablation: Introducing ablations in training on 4 LLMs
\end{enumerate}

For the first study, we choose model\_size as 384, 768, and 1024. \citet{mohapatra2025} studies a 768 size, but we experiment with lower and higher values as we build our own dataset, and it would be worth understanding the effects of different adaptor sizes across other datasets of different sizes than Gaddy and Klein.

For the ablation study, we introduce variations during training. Effects of changes such as cosine learning rate scheduling, gradient clipping, label smoothing, weight decay, data augmentation, layer normalisation, etc, are studied.

Table~\ref{tab:grid_vs_ablation} depicts the types of changes we make in our study and which parts of the pipeline are affected when. The parameters of the adaptor layers (the first three rows of Table~\ref{tab:grid_vs_ablation}) are varied in the grid search study, while everything else remains the same. For the ablation study, we keep the adaptor constant and vary ablation parameters (next 7 rows).

\begin{table}[ht]
\centering
\caption{Grid search vs.\ ablation: which parameters change in each study.}
\label{tab:grid_vs_ablation}
\begin{tabular}{lll}
\toprule
\textbf{Change} & \textbf{Grid search (model\_size)} & \textbf{Ablation (training recipe)} \\
\midrule
ResBlocks & Varies widths & Same (768/384/192) \\
BiLSTM & Varies dimensions & Same (192$\to$384) \\
Linear projection & Varies (MS//2$\to$llm\_hidden) & Same (384$\to$llm\_hidden) \\
Dropout & Fixed at 0.15 & Varies (0 vs 0.15) \\
LayerNorm & Fixed on & Varies (on vs off) \\
LR schedule & Fixed cosine & Varies (constant vs cosine) \\
Label smoothing & Fixed 0.1 & Varies (0 vs 0.1) \\
Augmentation & Fixed on & Varies (off vs on) \\
LLM mode & Fixed eval & Varies (train vs eval) \\
Prompt at inference & Fixed on & Varies (off vs on) \\
\bottomrule
\end{tabular}
\end{table}

\subsection{Stage 3: Frozen LLM Decoding}

This stage consist of a frozen LLM which maps the learned outputs of adaptor to words. The input embeddings are prepended by ``P1 = Unvoiced EMG:'' and appended by ``P2 = Prompt: Convert unvoiced EMG embeddings to text.'' These text embeddings are also converted to numbers using the LLM tokenizer having the mapping as $M : X_P \to H_P$, where $X_P = [P1, P2]$. On processing this through the LLM, we obtain the output embeddings for the EMG signals. These are the predicted logits $z_{t_s,y'}$ for each vocabulary item $y'$ at position $t_s$. Cross entropy loss is:

\begin{equation}
\mathcal{L} = -\sum_{t_s=1}^{T_s} \sum_{y' \in V} y'_{t_s} \log \left( \frac{\exp(z_{t_s,y'}/\tau)}{\sum_{v \in V} \exp(z_{t_s,v}/\tau)} \right),
\end{equation}

where $\tau$ is the temperature parameter and $y'_{t_s}$ is the true class at $t_s$. Figure~\ref{fig:pipeline_full} shows how the complete pipeline functions, where we concatenate embeddings of:

\begin{itemize}
    \item ``P1 = Unvoiced EMG:''
    \item EMG signals (multi-channel)
    \item ``P2 = Prompt: Convert unvoiced EMG embeddings to text.''
\end{itemize}

for the LLM to process. It outputs word embeddings which are then compared with the groundtruth word embeddings through cross-entropy loss, to output correct embeddings.

\begin{figure}[ht]
    \centering
    \includegraphics[width=\textwidth]{figures/image10.png}
    \caption{Trainable EMG adapter with frozen LLMs to transcribe text. Taken from \citet{mohapatra2025}.}
    \label{fig:pipeline_full}
\end{figure}

We experiment with training adaptors of model\_size = 768 for 8 different LLMs, including LLaMA-3.2-3B to establish a baseline for existing study \citet{mohapatra2025}. We also perform a grid search across 2 LLMs as a grid search for the best adaptor configuration.

Grid search is performed for the following LLMs:

\begin{enumerate}
    \item Gemma-2-2b-it
    \item Qwen2.5-3B
\end{enumerate}

Adaptors with model\_size = 768 are trained for the following LLMs:

\begin{enumerate}
    \item Gemma-2-2b-it
    \item Phi-3-mini-4k
    \item DeepSeek-R1-Distill-Qwen-1.5B
    \item LiquidAI/LFM2.5-1.2B-Instruct
    \item LLaMA-3.2-3B
    \item Qwen2.5-3B
    \item Stablelm-2-1\_6b-chat
    \item SmolLM2-1.7B-Instruct
\end{enumerate}

Apart from the two broad studies mentioned above, we perform two further studies that help us close the existing research gap, as follows:

\begin{enumerate}
    \item Channel drop from 8 to 4 and 3 EMG channels: We study how channel drop from 8 channels (research-grade, used by Gaddy) to 4 and 3 channels (consumer-grade, economically feasible) impacts performance.
    \item Custom dataset: While we experiment with train/test split on Gaddy dataset, we also collect our own data across two trials:
    \begin{enumerate}
        \item[i.] Same data as Gaddy (67-word date/time closed vocab)
        \item[ii.] A new 10-word (consisting of phonetically distinct fruit names) closed vocab
    \end{enumerate}
\end{enumerate}

Both studies use our best-performing LLM, the Gemma-2 model. We discuss them further in the next section.

\section{Results and Conclusions}

This section discusses all results of the experiments outlined in Section~\ref{sec:software_pipeline}.

\subsection{Grid Search: LLM and Adaptor Size}

First, we discuss the results obtained from the grid search on 2 LLMs in Table~\ref{tab:grid_search}. The `MS' column shows the model\_size evaluated for 3 values. `Hidden' column shows the llm\_hidden (LLM's input embedding space). Params represents the approximate number of parameters contained by the corresponding adaptor. WER is the Word Error Rate. CER is the Character Error Rate. We only run for 500 epochs; the last column shows the epoch when the best adaptor was obtained during training. It is worth noting that CER is important for the date/time dataset of Gaddy because small errors, such as a misprediction of `11 November' as `12 November', result in 50\% WER (as 11 and 12 are considered words by WER) but only 10\% CER (as individual characters are compared). Also, notably, running for more epochs may lead to better results.

\begin{table}[ht]
\centering
\caption{Grid search results across 2 LLMs and 3 model sizes.}
\label{tab:grid_search}
\small
\begin{tabular}{clcccccr}
\toprule
\textbf{Rank} & \textbf{LLM} & \textbf{MS} & \textbf{Hidden} & \textbf{Params} & \textbf{WER} & \textbf{CER} & \textbf{Best Epc} \\
\midrule
\rowcolor{gray!12} \textbf{1} & \textbf{gemma-2-2b-it} & \textbf{768} & \textbf{2304} & \textbf{6.0M} & \textbf{0.3285} & \textbf{0.1476} & \textbf{460} \\
2 & gemma-2-2b-it & 1024 & 2304 & 10.1M & 0.3577 & 0.2063 & 370 \\
3 & Qwen2.5-3B & 1024 & 2048 & 10.0M & 0.4088 & 0.1877 & 350 \\
4 & Qwen2.5-3B & 384 & 2048 & 1.8M & 0.4307 & 0.2049 & 290 \\
5 & gemma-2-2b-it & 384 & 2304 & 1.9M & 0.4453 & 0.2407 & 320 \\
6 & Qwen2.5-3B & 768 & 2048 & 5.9M & 0.4818 & 0.2550 & 250 \\
\bottomrule
\end{tabular}
\end{table}

We note that Google Gemma-2 2B model outperforms Qwen-2.5 3B for the 768 and 1024 MS. This result simply concludes that Gemma is more suited to be `fooled' into recognising and mapping learned EMG features to words. We conclude that Gemma has a better capacity to learn a broader context than simple word input embeddings. Gemma with MS=384 seems to have underfit the Gaddy dataset. Qwen model here shows unstable behavior where the MS=1024 variant is best performing, followed by MS=384 variant and MS=768. No definitive conclusion can be obtained from this.

Next, we demonstrate a thorough comparison between 8 LLMs, as shown in Table~\ref{tab:8llm_comparison}. Notably, this table only learnt adaptors for MS = 768 for every LLM and it is probable that the best performing LLM for each type be obtained using other adaptor sizes such as Qwen for MS = 1024 as shown in Table~\ref{tab:grid_search}. However, we believe MS = 768 encompasses the majority of LLMs at close-to-peak performance from random MS studies and from inspiration from \citet{mohapatra2025}.

\begin{table}[ht]
\centering
\caption{Comparison of 8 LLMs with model\_size = 768.}
\label{tab:8llm_comparison}
\resizebox{\textwidth}{!}{%
\begin{tabular}{clcccccr}
\toprule
\textbf{Rank} & \textbf{LLM} & \textbf{MS} & \textbf{Hidden} & \textbf{Params} & \textbf{WER} & \textbf{CER} & \textbf{Best Epc} \\
\midrule
\rowcolor{gray!12} \textbf{1} & \textbf{gemma-2-2b-it} & \textbf{768} & \textbf{2304} & \textbf{6.0M} & \textbf{0.3285} & \textbf{0.1476} & \textbf{460} \\
2 & Phi-3-mini-4k & 768 & 3072 & 6.3M & 0.3650 & 0.1777 & 390 \\
3 & DeepSeek-R1-Distill-Qwen-1.5B & 768 & 1536 & $\sim$5.5M & 0.4088 & 0.1748 & 310 \\
4 & LiquidAI/LFM2.5-1.2B-Instruct & 768 & 2048 & $\sim$5.9M & 0.4088 & 0.2020 & 300 \\
5 & LLaMA-3.2-3B & 768 & 3072 & 6.3M & 0.4307 & 0.2421 & 380 \\
6 & Qwen2.5-3B & 768 & 2048 & 5.9M & 0.4818 & 0.2550 & 250 \\
7 & stablelm-2-1\_6b-chat & 768 & 2048 & $\sim$5.9M & 0.5109 & 0.2908 & 440 \\
8 & SmolLM2-1.7B-Instruct & 768 & 2048 & $\sim$5.9M & 0.5401 & 0.3309 & 260 \\
\bottomrule
\end{tabular}%
}
\end{table}

We make the following conclusions from these results:

\begin{enumerate}
    \item \textbf{LLM Choice Matters Significantly}

The choice of LLM for the task of unvoiced speech recognition significantly affects the study's performance. The WER changes from 0.3285 to 0.5401 between Gemma and SmolLM2, a 64\% relative difference, even though the LLM is frozen. This signifies that the LLM's pre-training representation and quality heavily influence EMG decoding quality. Research into the documentation of LLM pre-training data, if possible, is hence encouraged. It is also noteworthy that LLM size is not directly proportional to performance, as the worst performing LLMs -- StableLM and SmolLM2 are the largest and one of the smallest LLMs in our study.

    \item \textbf{Hidden Dimension $\neq$ Performance}

Definitive conclusions cannot be made on the optimal llm\_hidden value. As we notice in Table~\ref{tab:hidden_dim}, DeepSeek has the smallest hidden dimension but ranks 3\textsuperscript{rd}, beating LLaMA-3.2-3B and Qwen2.5-3B, which have larger hidden dimensions and more parameters.

\begin{table}[ht]
\centering
\caption{Hidden dimension vs.\ performance comparison.}
\label{tab:hidden_dim}
\begin{tabular}{ccl}
\toprule
\textbf{Hidden} & \textbf{Best WER} & \textbf{LLM(s)} \\
\midrule
\rowcolor{gray!12} 2304 & \textbf{0.3285} & Gemma-2 \\
3072 & 0.3650 & Phi-3 \\
\rowcolor{gray!12} 1536 & 0.4088 & DeepSeek \\
2048 & 0.4088--0.5401 & LiquidAI, LLaMA, Qwen, StableLM, SmolLM \\
\bottomrule
\end{tabular}
\end{table}

    \item \textbf{Convergence Speed Variation}

There is no definitive correlation between convergence speed and final performance, as we note that DeepSeek and Qwen converge fastest around 250-310 epochs, whereas StableLM performs poorly despite converging at 440 epochs.
\end{enumerate}

\subsection{Ablation Study: Training Recipe}

Next, we discuss results from ablation study performed on 4 LLMs, as follows:

\begin{enumerate}
    \item gemma-2-2b-it
    \item LLaMA-3.2-3B-Instruct
    \item Phi-3-mini-4k-instruct
    \item Qwen2.5-3B-Instruct
\end{enumerate}

In this study, we find that each adaptor responds slightly differently to an introduced ablation in training.

Table~\ref{tab:ablation} displays the impact of ablations on Gemma-2. We discard exact results for other LLMs and instead show ablation-wise results for all LLMs in the last column of the table. We select the best model from S6 ablation (6 stacked ablations) as further ablations hurt Gemma's performance.

\begin{table}[ht]
\centering
\caption{Ablation study results on Gemma-2 and summary across 4 LLMs.}
\label{tab:ablation}
\resizebox{\textwidth}{!}{%
\begin{tabular}{clcccc}
\toprule
\textbf{Stage} & \textbf{Improvement Added} & \textbf{WER} & \textbf{CER} & \textbf{$\Delta$WER} & \textbf{Cross-LLM (4 LLMs)} \\
\midrule
S0 & Baseline (all OFF) & 0.9635 & 0.7550 & --- & --- \\
S1 & + LLM eval mode & 0.9635 & 0.7550 & 0.0000 & 0/4 helped \\
\rowcolor{gray!12} S2 & + Prompt-consistent inference & 0.6277 & 0.3868 & $-$0.3358 & 4/4 helped \\
S3 & + Cosine LR schedule & 0.6131 & 0.4241 & $-$0.0146 & 2/4 helped \\
\rowcolor{gray!12} S4 & + Gradient clipping & 0.5839 & 0.3596 & $-$0.0292 & 3/4 helped \\
\rowcolor{gray!12} S5 & + Adaptor dropout (0.15) & 0.4307 & 0.2249 & $-$0.1532 & 4/4 helped \\
\rowcolor{gray!12} \textbf{S6} & \textbf{+ Output LayerNorm} & \textbf{0.3212} & \textbf{0.1447} & $\mathbf{-}$\textbf{0.1095} & \textbf{3/4 helped} \\
S7 & + Label smoothing (0.1) & 0.3285 & 0.1590 & $+$0.0073 & 2/4 helped \\
S8 & + Data augmentation & 0.3723 & 0.1777 & $+$0.0438 & 1/4 helped \\
S9 & + Weight decay (0.05) & 0.3285 & 0.1476 & $-$0.0438 & 2/4 helped \\
\bottomrule
\end{tabular}%
}
\end{table}

The following are our conclusions:

\begin{enumerate}
    \item \textbf{Prompt-Consistent Inference:}

This is a bug fix as the original codebase has a train/test mismatch. This also justifies our WER of 0.43 vs 0.49 in the base paper for the LLaMA run. This ablation helps all 4 LLMs we study in improving WER.

    \item \textbf{Output LayerNorm:}

This ablation improves WER for 3 out of 4 LLMs. This is attributed to the fact that adaptor outputs raw linear projections into llm\_hidden space, but LayerNorm is required to convert them into LLM's specific scale, for better processing.

    \item \textbf{Regularisation:}
    \begin{enumerate}
        \item[i.] Dropout (S5): helped all 4 LLMs, universally beneficial with only 500 training samples
        \item[ii.] Grad clipping (S4): helped 3 of 4 LLMs, stabilises training
        \item[iii.] Label smoothing (S7): helped and hurt 2 LLMs each. Huge for Phi-3 ($-0.10$) but slightly hurt Qwen ($+0.007$)
        \item[iv.] Weight decay (S9): again, helps and hurts 2 LLMs each, with no specific correlation
    \end{enumerate}

    \item \textbf{Data augmentation:}

Not the best ablation, as it only aids Qwen in increasing WER. We believe data augmentation adds noise that confuses our models during training for the 67-phrase closed vocab dataset.

    \item \textbf{Best Stage:}

This varies by LLM and implies that the optimal training recipe is LLM-specific. The best stages are as follows:
    \begin{enumerate}
        \item[i.] Gemma-2: S6
        \item[ii.] Phi-3: S7
        \item[iii.] Qwen: S8
        \item[iv.] LLaMA: S7
    \end{enumerate}
\end{enumerate}

Overall, we conclude with certainty that the majority of ablations helped us improve our adaptors' performance on the task of unvoiced speech recognition.

\subsection{Channel Reduction Analysis}

In this section, we discuss results from the first further study mentioned in Section~\ref{sec:software_pipeline}.

The selection of the best 4 and 3 channels is done using three independent metrics and combining the ranks via Borda count. The methods and the quantities they measure are stated in Table~\ref{tab:channel_methods}. Each method produces a ranking 1--8. The Borda score is variance\_rank + uniqueness\_rank + energy\_rank. The top $N$ channels by lowest score are selected.

\begin{table}[ht]
\centering
\caption{Channel selection methods and the quantities they measure.}
\label{tab:channel_methods}
\begin{tabular}{lp{5.5cm}p{5.5cm}}
\toprule
\textbf{Method} & \textbf{What it measures} & \textbf{Rationale} \\
\midrule
Variance & Total feature variance per channel & Higher variance = more informative signal \\[4pt]
Uniqueness & $1 - \text{mean correlation with other channels}$ & Lower correlation = non-redundant info \\[4pt]
Energy & L2 norm of channel features & Stronger muscle activation = better SNR \\
\bottomrule
\end{tabular}
\end{table}

Table~\ref{tab:channel_results} displays results when we perform channel reduction on the Gaddy dataset. The first row is the standard Gemma-2 results which have been obtained earlier. The second row corresponds to results with 4 channels and the third row to 3 channels.

Here are our conclusions from the results:

\begin{enumerate}
    \item \textbf{Linear degradation:} We see a drop of 0.08 WER and 0.07 WER approximately as we reduce the number of channels. This implies that based on budget and accuracy constraints, one can deduce the correct number of channels to be used.

    \item \textbf{Convergence rate:} We notice faster converge (reduction in number of epochs) as we reduce channels. This explains the fact well that less data (due to fewer features) gets learned by the model faster, although it converges to a worse optimum.

    \item \textbf{Establish state-of-the-art results:} Our worst-performing 3-channel Gemma-2 model itself beats \citet{mohapatra2025} results by approximately 0.02 WER, using the LLaMA-3.2-3B model with 8 channels. This is a significant improvement, especially considering resource-constrained train/test environments.
\end{enumerate}

\begin{table}[ht]
\centering
\caption{Channel reduction results on the Gaddy dataset using Gemma-2-2b-it (MS=768, Hidden=2304). The ResBlocks column shows the input dimension change; LSTM and Projection remain fixed.}
\label{tab:channel_results}
\begin{tabular}{cccccccc}
\toprule
\textbf{Channels} & \textbf{Features} & \textbf{Params} & \textbf{WER} & \textbf{CER} & \textbf{Best Epc} & \textbf{ResBlocks input} \\
\midrule
\rowcolor{gray!12} 8 & 112 & 6.0M & \textbf{0.3285} & \textbf{0.1476} & 460 & $112 \to 768/384/192$ \\
4 & 56 & $\sim$5.97M & 0.4088 & 0.2077 & 360 & $56 \to 768/384/192$ \\
3 & 42 & $\sim$5.95M & 0.4745 & 0.2507 & 290 & $42 \to 768/384/192$ \\
\bottomrule
\end{tabular}
\end{table}

\subsection{Personal EMG Dataset Experiments}

\subsubsection{Custom dataset -- similar to Gaddy}

We collected this data for 438 samples. We use three methods:

\begin{enumerate}
    \item Zero-shot: Inference only using the adaptor trained on the Gaddy dataset
    \item Fine-tune: Train a learned adaptor (from Gaddy) further on our dataset
    \item From-scratch: Train an adaptor from scratch on only our dataset
\end{enumerate}

The results are displayed in Table~\ref{tab:custom_gaddy}. It is evident that from-scratch training on only our dataset yields much better results than otherwise. This can be attributed to \citet{mohapatra2025}, which states that EMG signals are highly person-specific and that there is significant variation in how muscles move between people. Also, Gaddy used 8 channels while we use only 4, and the sensor quality and placement vary highly between the two studies, leading to sub-optimal fine-tuned results. It is also noteable that zero-shot transfer completely fails and the Gaddy-trained adaptor produces random output on personal EMG data.

\begin{table}[ht]
\centering
\caption{Custom dataset results (Gaddy-style 67-word vocabulary).}
\label{tab:custom_gaddy}
\begin{tabular}{lcc}
\toprule
\textbf{Experiment} & \textbf{Best WER} & \textbf{Best CER} \\
\midrule
A) Zero-shot (Gaddy weights $\to$ our data) & 1.032 & 0.931 \\
B) Fine-tune (Gaddy weights $\to$ train on our data) & 0.925 & 0.804 \\
\rowcolor{gray!12} \textbf{C) From scratch (random init $\to$ our data)} & \textbf{0.871} & \textbf{0.726} \\
\bottomrule
\end{tabular}
\end{table}

\subsubsection{Custom dataset -- New 10-word dataset}

10 words which are phonetically distinct fruits and vegetables are used in this study. We collect 250 samples, 25 samples per class, ensuring a dataset with balanced classes. From another study, we find that the optimum model\_size for our small dataset is 384. This is less in comparison to 768, which worked well for a larger dataset such as Gaddy's. Hence, we present results in Table~\ref{tab:custom_10word} wherein we only train/test adaptor with model\_size = 384 on our dataset. This is also conducted using our best-performing Gemma-2 LLM.

\begin{table}[ht]
\centering
\caption{Custom 10-word dataset results with varying channel counts.}
\label{tab:custom_10word}
\begin{tabular}{cccc}
\toprule
\textbf{Channels} & \textbf{Features} & \textbf{WER} & \textbf{Word Accuracy} \\
\midrule
1ch & 14 & 0.3600 & 64\% \\
2ch & 28 & 0.2400 & 76\% \\
\rowcolor{gray!12} \textbf{3ch} & \textbf{42} & \textbf{0.0800} & \textbf{92\%} \\
\rowcolor{gray!12} \textbf{4ch} & \textbf{56} & \textbf{0.0800} & \textbf{92\%} \\
\bottomrule
\end{tabular}
\end{table}

This table is significant as we make the following conclusions:

\begin{enumerate}
    \item \textbf{Justifies our whole study:} Consumer-grade EMG setup works even with just 3 channels and 250 training samples on a closed 10-word vocabulary. We prove that this is especially helpful to set up a simple setup for people who have completely lost their voice and would like to communicate very simple words for basic functions.

    \item \textbf{Channel-wise results:} A single channel achieves 64\% accuracy when targeting the muscles we highlighted previously. We see no improvement from 3 channels to 4 channels due to a combination of 2 factors: possible muscle overlap, and noisy 4\textsuperscript{th} channel signals, which impact its performance.

    \item \textbf{Economic feasibility:} Once again, we demonstrate that, given resource constraints, the appropriate number of channels can be identified and used.
\end{enumerate}

\subsubsection{Discussion of Personal Dataset Results}

While the results above are encouraging, several practical challenges emerged during our personal data experiments that warrant discussion.

\paragraph{Trial 1: Complexity vs.\ Data.} The 67-word Gaddy-style vocabulary proved too complex relative to our 438-sample dataset. With 426 unique phrases from only 82 unique tokens, many phrase combinations appeared only once, leaving the adaptor with insufficient repetition to learn reliable EMG-to-text mappings. Additionally, our attempt to fuse EMG samples from different participants (collected across separate sessions) failed to improve performance. This is consistent with prior findings that sEMG signals are highly person-specific \citep{mohapatra2025}---differences in facial anatomy, muscle activation patterns, and electrode impedance between individuals render cross-subject data fusion ineffective without explicit domain adaptation.

\paragraph{Trial 2: The Inference Gap.} Although the 10-word fruit classification achieved 92\% accuracy on the development set, we observed a noticeable drop during standalone inference sessions conducted after the main data collection. We attribute this to three factors:
\begin{itemize}
    \item \textbf{Placement sensitivity:} Even a 2--3\,mm shift in sensor placement between sessions produced measurably different signals, owing to the high density of overlapping facial muscles in the perioral region.
    \item \textbf{Skin integrity:} Repeated electrode application caused mild skin irritation, particularly on the neck, which degraded the electrode--skin contact impedance and introduced noise in later sessions.
    \item \textbf{Hardware variance:} While the mapping from muscle location to Arduino analog pin was kept consistent, the physical sensor-to-location assignment was not always identical across sessions, introducing subtle gain and offset differences.
\end{itemize}

These observations underscore that robust real-world deployment of consumer-grade EMG silent speech systems requires careful attention to electrode placement reproducibility, session-to-session calibration, and potentially online adaptation techniques.

\section{Project Novelties}

The authors believe that they have made significant research contributions from the results obtained, alongside introducing several novelties, as listed below. Existing literature does not cover the following aspects:

\begin{enumerate}
    \item \textbf{Comparison of 8 LLMs:} We cover a wide range of LLMs with different capabilities and pre-training strategies (Table~\ref{tab:8llm_comparison}). For every LLM, we train an adaptor. \citet{mohapatra2025} only covers LLaMA-2-7B and LLaMA-3.2-3B and achieves a WER of 0.49. We improve this to 0.33 using Gemma-2. Prior EMG-to-text literature has not examined how frozen LLMs affect silent speech decoding.

    \item \textbf{Ablation study:} We holistically evaluate training across 9 incremental ablations across 4 LLMs, showing which generalise and which are LLM-specific (Table~\ref{tab:ablation}).

    \item \textbf{Consumer-grade EMG Hardware Validation:} Previous studies either use the Gaddy dataset only and do not perform custom dataset inference, or use research-grade electrodes. We make a sharp shift and collect both Gaddy-style data to test benchmarks and our own phonetically distinct data to establish new baselines.

    \item \textbf{Channel reduction analysis:} No prior research work systematically studies the reduction in EMG channels. We develop a principled pipeline to select a number of channels based on variance, uniqueness, and energy of signals and validate it with data-loss metrics such as PCA, linear reconstructions, and cross-correlation (Tables~\ref{tab:channel_methods} and \ref{tab:channel_results}).
\end{enumerate}

Based on this, we aim to publish a research paper at an IEEE/ACM conference. UIST (ACM Symp.\ on User Interface Software and Tech) and ACM Multimedia (ACM MM) are some of the conferences we are targetting, with work on research paper writing having already been started.

\section{Future Work}

Several directions emerge naturally from this project and could strengthen both the system's practical utility and its research contributions.

\paragraph{Larger and More Diverse Datasets.} Our personal dataset experiments were limited by sample size (438 and 757 utterances respectively) and number of participants. Collecting data from a larger cohort with multiple recording sessions per participant would enable a more rigorous evaluation of cross-session and cross-subject generalisation. Expanding the vocabulary beyond the closed 67-word set toward open-vocabulary decoding---potentially by leveraging the LLM's full generative capacity---remains a key long-term goal.

\paragraph{Session-to-Session Calibration.} The inference gap observed in Trial~2 highlights the need for online adaptation techniques that can compensate for electrode placement drift between sessions. Approaches such as few-shot calibration, where a short re-recording session updates the adaptor's input normalisation, or test-time adaptation methods could significantly improve deployment robustness.

\paragraph{Extended LLM and Adaptor Exploration.} While we compared 8 frozen LLMs, fine-tuning strategies such as LoRA \citep{mohapatra2025} were not explored in depth. Additionally, our ablation study was limited to 500 training epochs; extending training and applying learning rate warm-up or cyclical schedules may yield further improvements, particularly for the larger adaptor configurations.

\paragraph{Channel Dropout Pre-training.} Inspired by \citet{hwang2026}, applying channel dropout during training on the full 8-channel Gaddy dataset before fine-tuning on 4 or 3 channels could improve reduced-channel performance beyond our current results, which train directly on the subset.

\paragraph{Real-time Deployment.} Although we demonstrated a working Arduino-based data acquisition pipeline, end-to-end real-time inference---from EMG capture through feature extraction to LLM decoding---has not yet been optimised for latency. Deploying a quantised adaptor on edge hardware while streaming to a server-hosted LLM would be a practical next step toward a usable silent speech device.

\paragraph{Publication.} We plan to prepare a 6--7 page paper for submission to an IEEE or ACM venue, consolidating the LLM comparison, ablation study, channel reduction analysis, and personal dataset results into a concise research contribution.

\clearpage
% Harvard references (agsm style)
\bibliographystyle{agsm}
\bibliography{ref}
\end{document}